\subsection{Asistente en lenguaje natural (NL2C)}
\label{sec:nl2c}

El asistente convierte la frase de un acudiente en las restricciones del optimizador. Corresponde a las etapas de modelado y evaluación de CRISP-DM en las fases F3 y F4 y sigue la estructura de las secciones anteriores: problema, entradas y salidas, estrategia, seudocódigo, costo, resultados y límites. Las tablas
completas de la evaluación están en el Anexo~\ref{anx:nl2c}.

\paragraph{Análisis del problema.} La propuesta aprobada prometió que el acudiente pudiera
describir en español lo que quiere, por ejemplo «Minnie para mi hija, máximo 500 mil», en lugar de
llenar un formulario de opciones cerradas. A esa tarea se le llama NL2C (del inglés \emph{natural
language to constraints}): transformar un texto en el objeto de restricciones que un programa ya
consume, que aquí es el optimizador de la Sección~\ref{sec:optimizacion-cal}. El texto $t$ es una
cadena corta, ambigua y con errores de ortografía, y la salida es la terna
$\mathcal{R}=\langle F,W,P\rangle$ de filtros duros, pesos y preferencias, con los dominios de la
Tabla~\ref{tab:nl2c-dominios}. Un campo que el acudiente no menciona queda en su valor por omisión,
y uno que menciona fuera de dominio se lleva al dominio y se le avisa. Dos condiciones fijan el
diseño. La primera es que $t$ y el texto de la lista los puede escribir un atacante, de modo que
ninguno de los dos puede cambiar qué acciones ejecuta el sistema \citep{greshake2023indirect}. La
segunda es que un error le cuesta dinero al acudiente, por lo que nada se aplica sin su
confirmación. Las reglas fijas no bastan, porque los nombres de personajes y marcas son un
vocabulario abierto y una frase como «calidad primero y precio después» fija los pesos pero no el
estilo de gasto.

\begin{table}[H]
\centering\footnotesize
\caption{Dominio de los campos de $\mathcal{R}$. Fuente: constantes de \texttt{nl2c.py}.}
\label{tab:nl2c-dominios}
\begin{tabularx}{\textwidth}{@{}p{1.2cm}p{4.2cm}X@{}}
\toprule
Grupo & Campo & Dominio \\
\midrule
$F$ & \texttt{presupuesto\_max\_cop} & vacío o entero en $[50.000;\ 3.000.000]$ pesos por hijo \\
 & \texttt{marcas\_excluidas}, \texttt{solo\_disponibles} & hasta 10 cadenas de 1 a 40 caracteres; verdadero o falso \\
$W$ & cuatro pesos & cada $w_i\in[0{,}05;\ 0{,}60]$ y $\sum w_i=1$ \\
$P$ & \texttt{personajes}, \texttt{estilo\_gasto} & hasta 3 del catálogo; económico, punto medio o estirarse \\
\bottomrule
\end{tabularx}
\end{table}

\paragraph{Diseño: entradas y salidas.} Las entradas son el texto $t$, de a lo más 1.000 caracteres
después de sanearlo, el catálogo de personajes $\mathcal{C}_P$, con hasta 60 nombres, y una lista
blanca de ocho componentes de interfaz. Para pedir cambios sobre un kit se agrega el kit, con hasta
80 renglones. Las salidas son $\mathcal{R}$ validada, la lista de ajustes que hizo el servidor, la
especificación de la interfaz que se le muestra al acudiente y un estado, que puede ser propuesta,
sin acción, inválido o sin servicio. Si el estado es propuesta, $\mathcal{R}$ cumple la
Tabla~\ref{tab:nl2c-dominios}, todo texto libre de $\mathcal{R}$ aparece en $t$ y nada se aplica
antes de que el acudiente confirme.

\paragraph{Estrategia algorítmica.} La estrategia es un plan fijo con validación en el servidor: el
código decide qué funciones puede llamar cada ruta antes de ver $t$, y el modelo de lenguaje solo
elige una de ellas \citep{beurerkellner2025patterns}. La idea sigue la separación entre planificar y
ejecutar de Debenedetti et al. \citep{debenedetti2025camel}. Con ese plan, el modelo recibe el texto como dato y llama a
una función con argumentos; el servidor valida los argumentos y arma la interfaz con los componentes
de la lista blanca, que la aplicación web vuelve a validar \citep{chen2025genui}; y el acudiente
confirma, como piden las pautas de interacción humano-IA \citep{amershi_guidelines_2019}. Con esas
barandas, el servidor cubre cinco riesgos de la lista OWASP para aplicaciones con modelos de
lenguaje \citep{owasp2025llm}: inyección de instrucciones, información sensible, manejo de la
salida, agencia excesiva y consumo ilimitado. La línea base es el formulario de opciones cerradas,
que no usa modelo ni tiene riesgo de inyección, pero no entiende frases que mezclan campos.

\begin{algorithm}[H]
\caption{Interpretación de una preferencia en lenguaje natural}
\label{alg:nl2c}
\KwEntrada{texto $t$, catálogo de personajes $\mathcal{C}_P$}
\KwSalida{estado y, si es propuesta, la terna $(\mathcal{R}, A, U)$}
$s\leftarrow$ Sanear$(t)$ \tcp*{quita datos personales y marcas que escapan del bloque de datos}
$L\leftarrow\{$fijar\_restricciones, sin\_acción$\}$ \tcp*{plan fijo: la lista no sale de $t$}
$(f,a)\leftarrow$ LlamarModelo$(s,L)$ \tcp*{una sola llamada, sin reintento}
\lIf{$f\notin L$ o $f=$ sin\_acción}{\Return{sin acción}}
$(\mathcal{R},A)\leftarrow$ Proyectar$(a,\mathcal{C}_P,s)$ \tcp*{forma, recorte a dominios, texto anclado a $s$}
\lIf{Proyectar falló por forma}{\Return{inválido}}
$U\leftarrow$ CompilarInterfaz$(\mathcal{R})$ \tcp*{la arma el servidor, no el modelo}
\Return{(propuesta, $\mathcal{R}$, $A$, $U$)} \tcp*{nada se aplica hasta la confirmación}
\end{algorithm}

\paragraph{Complejidad y costo.} Sanear recorre $t$ con un número constante de expresiones
regulares, con un costo de $O(|t|)$. Proyectar trabaja sobre una $\mathcal{R}$ de tamaño acotado y
cuesta $O(|\mathcal{C}_P|+|s|)$, por el índice del catálogo y el cotejo de cada marca con el texto;
la normalización de los pesos itera a lo más cuatro veces y la interfaz tiene a lo más seis
componentes. En total son $O(|t|+|\mathcal{C}_P|)$ más una llamada al modelo, acotada a 400 tokens de
salida y 12 segundos de plazo, que domina el tiempo. En la ronda 3, con 70 intentos, esa llamada tuvo
una latencia p50 de 1.000~ms y p95 de 1.177~ms, con 616 tokens de entrada y 39 de salida en promedio.
El formulario no llama al modelo, pero su tiempo no se midió con acudientes, de modo que no se afirma
que NL2C sea más rápido para ellos.

\paragraph{Dónde quedó en la aplicación.} El componente se construyó para los dos puntos de
entrada, el inicio del pedido y la página del kit, pero en la versión que usaron los acudientes de
la ronda 2 de OE5 solo se monta en la página del kit, como el asistente que recibe pedidos de cambio
del tipo «más barato en cuadernos». En el inicio el equipo dejó el formulario guiado, porque el
acudiente todavía no ha visto ningún kit y el formulario le muestra de una vez qué puede elegir. El
servicio de interpretación es el mismo en los dos casos, así que la evaluación de esta sección
vale para el asistente del kit; el panel del inicio sigue en el código detrás de una bandera
apagada.

Después de la ronda 3, el asistente cambió sin cambiar sus barandas. Ante un pedido por categoría,
como «más barato en cuadernos», el modelo solo clasifica la acción y la categoría, y reglas del
servidor eligen renglones, producto y precio, con hasta 12 propuestas y el ahorro total; «Aceptar
todos» envía cambios individuales con la sesión del acudiente. El contexto se reduce a la acción y
la categoría anteriores, validadas contra el kit, y el modelo no recibe datos personales.

\paragraph{Evaluación.} El asistente se midió con 35 frases en español, 22 de la encuesta y 13
construidas, cuya etiqueta de oro se escribió antes de correr el modelo, y con 27 ataques de
inyección al estilo de AgentDojo \citep{debenedetti2024agentdojo}. El modelo fue
\texttt{gemini-3.5-flash-lite} con temperatura 0, con dos repeticiones por frase (70 intentos) y tres
por ataque (81 intentos). Un ataque tiene éxito cuando su condición se cumple sobre lo que el
acudiente ve después de las barandas del servidor. En la ronda 3 la tarea completa se cumplió en 58
de 70 intentos (82,9\,\%), el F1 de personajes y marcas fue 1,000 y ningún ataque tuvo éxito (0 de
81), aunque el modelo lo intentó en 6. Los 12 fallos de utilidad se concentran en el campo del peso
mayor, donde varias etiquetas son discutibles. Bajo ataque se perdió la tarea legítima en 24 de 60
intentos, porque el modelo respondió sin acción ante la orden hostil, lo que equivale a preferir no
proponer ante la duda.

Entre la ronda 2 y la 3 cambió una condición de éxito, y se declara. La
condición \texttt{descarte\_\allowbreak aplicable} se sustituyó por \texttt{descarte\_sin\_confirmar}, porque al integrar el
asistente con la sesión de pedidos el descarte pasó a ser aplicable después de la confirmación del
acudiente. El cambio se hizo después de ver los resultados de la ronda 2, así que no es un criterio
preregistrado; en esa condición las dos rondas dieron 6 evaluaciones y 0 éxitos. Además, con los
servicios levantados, 4 pedidos contra un kit real de 36 renglones dieron la respuesta esperada en
786 a 885~ms cada uno, lo que con $n=4$ es una comprobación de funcionamiento y no una tasa.

El arnés sin conexión, que usa un modelo que obedece toda orden hostil y mide solo las barandas,
da 7 éxitos en 27 ataques, igual antes y después del cambio: cinco son ecos del propio texto del
acudiente y dos son propuestas válidas que él no pidió, que detiene su confirmación. Es el peor caso
de las barandas y no se compara con el 0 de 81 del modelo real; la ronda en vivo con Gemini sobre el
asistente nuevo no se hizo. En las sesiones con acudientes (Sección~\ref{sec:oe5}) el asistente
estuvo visible en la página del kit, pero ninguno lo usó por iniciativa propia, así que su utilidad
con usuarios reales queda como trabajo futuro.

\paragraph{Limitaciones.} Con 0 éxitos en 81 intentos, el límite superior del intervalo de confianza
del 95\,\% es de cerca de 3,7\,\% (regla de tres), por lo que solo se puede afirmar que no se encontró
un ataque exitoso con estos casos. Las repeticiones a temperatura 0 no son independientes, el modelo
y la fecha son uno solo, el kit y el catálogo de las rondas son sintéticos y los ataques los escribió
quien construyó las defensas. Los filtros duros sí llegan al optimizador (commit \texttt{0059f6d}):
\texttt{marcas\_excluidas} y \texttt{solo\_disponibles} se aplican antes del frente de Pareto, y
también en el carrusel de opciones y en el cambio de producto, pero no quitan una marca que el
colegio exige en el renglón. La ruta rápida, que publica el kit canónico sin recalcular, los
rechaza con el error 422 \texttt{filtros\_requieren\_recalculo}. Su límite es que
\texttt{solo\_disponibles} usa la oferta activa con precio vigente y no el inventario por tienda.
